\section*{الفصل \ref{ch:b2:cell-signalling} --- الرسل الكيميائية ونقل الإشارة}

\begin{solution}{exo:b2:cell-signalling:1}
الأنسولين: صماوي، ذائب في الماء (ببتيد). والكورتيزول: صماوي،
ذائب في الدسم. والأستيل كولين عند مشبك: مشبكي، ذائب في الماء.
والهيستامين عند جرح: نظير صماوي، ذائب في الماء. والإستراديول:
صماوي، ذائب في الدسم. والأدرينالين من الكظر: صماوي،
ذائب في الماء. وأكسيد النتريك: نظير صماوي، ذائب في الدسم (غاز
يعبر الأغشية).
\end{solution}

\begin{solution}{exo:b2:cell-signalling:2}
يرتبط الأدرينالين بالمستقبل (والإطفاء: التفكك، وابتلاع
المستقبل)؛ ويفعّل المستقبل \omterm{prop:b2:cell-signalling:gpcr}{بروتين G}، وGTP بدل GDP (والإطفاء:
حلمهة GTP في ثوانٍ)؛ ويفعّل \omterm{prop:b2:cell-signalling:gpcr}{بروتين G} أدنيلات سيكلاز،
وهي تصنع \omterm{prop:b2:cell-signalling:gpcr}{AMP الحلقي} (والإطفاء: الفوسفودي إستيراز)؛ ويفعّل \omterm{prop:b2:cell-signalling:gpcr}{AMP الحلقي} كيناز A
(والإطفاء: زوال \omterm{prop:b2:cell-signalling:gpcr}{AMP الحلقي})؛ ويفسفر كيناز A كيناز الفسفوريلاز،
الذي يفسفر فسفوريلاز الغليكوجين (والإطفاء: الفوسفاتازات)؛
ويطلق الفسفوريلاز غلوكوز-1-فوسفات من الغليكوجين، فيُحوَّل
ويُصدَّر غلوكوزًا.
\end{solution}

\begin{solution}{exo:b2:cell-signalling:3}
ترتبط المستقبلات السطحية بالرسل الذائبة في الماء عند الغشاء
وتفعل في ثوانٍ عبر الرسل الثواني والكينازات، فتغير
فعالية بروتينات موجودة. أما \omterm{prop:b2:cell-signalling:nuclear}{المستقبلات النووية} فترتبط بالرسل الذائبة
في الدسم داخل الخلية وتفعل في ساعات \omterm{prop:b2:cell-differentiation:transcriptional}{عوامل
استنساخ}، فتغير أي البروتينات يُصنع.
\end{solution}

\begin{solution}{exo:b2:cell-signalling:4}
\omterm{prop:b2:cell-signalling:gpcr}{AMP الحلقي}: تصنعه أدنيلات سيكلاز، ويتلفه
الفوسفودي إستيراز. والكالسيوم: يدخل عبر قنوات أو يغادر
الشبكة الهيولية الباطنة جوابًا لثلاثي فوسفات الإينوزيتول، وتزيله
المضخات. وثلاثي فوسفات الإينوزيتول (مع ثنائي أسيل الغليسرول): يصنعه
الفوسفوليباز C من دسم غشائي، وتزيله الفوسفاتازات.
\end{solution}

\begin{solution}{exo:b2:cell-signalling:5}
$\theta = [L]/(K_d + [L])$: أي $0.09$ و$0.5$ و$0.91$ و$0.99$. ومضاعفة
المستقبلات تترك الكسر على حاله وتضاعف عدد
المشغول --- ومن ثم الاستجابة.
\end{solution}

\begin{solution}{exo:b2:cell-signalling:6}
$50\times 200\times 100\times 300 = 3\times 10^{8}$. ومئة مستقبل
مشغول تعطي $3\times 10^{10}$ جزيء ناتج في الثانية --- وهو
حد أعلى، لأن المراحل تتشبع.
\end{solution}

\begin{solution}{exo:b2:cell-signalling:7}
يبقي \omterm{prop:b2:cell-signalling:gpcr}{بروتين G} المقفل السيكلاز مشتغلًا، فيبقى \omterm{prop:b2:cell-signalling:gpcr}{AMP الحلقي} مرتفعًا، ويبقي كيناز A
قناة الكلور في خلايا الأمعاء مفسفرة
ومفتوحة، فيُفرز الكلور، ويتبعه الصوديوم والماء إلى
اللمعة --- لترات في اليوم. والقفل تعديل تساهمي في \omterm{prop:b2:cell-signalling:gpcr}{بروتين
G}، لا يُنقض إلا حين تُستبدل الخلايا نفسها بعد
أيام.
\end{solution}

\begin{solution}{exo:b2:cell-signalling:8}
$0.9\times 10^{-6}\,\text{مول/لتر}\times 2\times 10^{-12}\,\text{لتر} =
1.8\times 10^{-18}$ mol، أي $1.1\times 10^{6}$ شاردة (وأكثر في الواقع، لأن
الدارئات تربط معظم ما يدخل). وقناة واحدة عند $10^{6}$ في الثانية
تستغرق ثانية؛ وألف قناة تفعلها في مليثانية،
وهو المقياس الزمني الذي تحتاجه إشارة.
\end{solution}

\begin{solution}{exo:b2:cell-signalling:9}
يلتقي الطريقان على الإنزيمات نفسها: فكيناز A التابع للغلوكاغون
يفسفر مصنعة الغليكوجين (فتُطفأ) وكيناز الفسفوريلاز
(فيشتغل)؛ وشلال الأنسولين يفعّل الفوسفاتاز التي تنزع
الفوسفات. وحالة الإنزيمات يقررها ميزان فعالية الكيناز
والفوسفاتاز، فالمهم هو نسبة الهرمونين
أحدهما إلى الآخر؛ وحين يرتفعان معًا يغلب الفوسفاتاز عادةً فيُخزن
الغليكوجين.
\end{solution}

\begin{solution}{exo:b2:cell-signalling:10}
الأدرينالين: ينتشر \omterm{prop:b2:cell-signalling:gpcr}{AMP الحلقي} عبر خلية قطرها \qty{20}{\micro m} في أقل من
ثانية ($x^{2}/2D$ مع $D \approx \qty{3e-10}{m^2/s}$)، وكل
إنزيم في الشلال موجود سلفًا --- فثوانٍ. أما الكورتيزول: فالنسخة
تستغرق \qty{60}{s}، والبروتين \qty{100}{s}؛ ومن بضع
عشرات من mRNA يحمل كل منها عشرة ريباسات أو نحوها تصنع الخلية بضع
مئات من جزيئات الإنزيم في الدقيقة، فالآلاف اللازمة لأثر
قابل للقياس تستغرق ساعات، بعد الدقائق التي يحتاجها الكورتيزول
ليدخل ويبلغ DNA.
\end{solution}

\begin{solution}{exo:b2:cell-signalling:11}
يشير المستقبل باستمرار: فشلال كيناز MAP يدفع
الخلية إلى الانقسام دون أي عامل نمو، وهي تتجاهل
غياب الإشارة التي تكبحها عادةً. ولأن \omterm{def:b2:genome-diversification:mutation}{طفرة} واحدة
تزيل ضابطًا على الانقسام، كانت هذه المستقبلات (والكينازات
التي دونها) من أشيع المورّثات الورمية. والدواء الذي
يشغل موقع ATP في الكيناز يوقف الفسفرات
ويُخرس المستقبل --- وهو مبدأ عدة أدوية للسرطان.
\end{solution}

\begin{solution}{exo:b2:cell-signalling:12}
الصماوية: تبلغ كل خلية في دقيقة، ورخيصة لكل خلية، ولا تستطيع
مخاطبة خلية بعينها، وتفعل من دقائق إلى أيام --- ولا غنى عنها
لحالات الجسم كله (الصيام، والنمو، والتكاثر). والعصبية:
مليثانية، وهدف واحد، وأسلاك مكلفة، وتفعل مليثوان
--- ولا غنى عنها للحركة والمنعكسات السريعة. ويُستعمل الاثنان
للرسالة نفسها حيث تُطلب السرعة والمدى معًا: فالأعصاب
الودّية ولب الكظر كلاهما يسلّم
النورأدرينالين أو الأدرينالين عند فزع.
\end{solution}

\begin{solution}{pb:b2:cell-signalling:1}
\textbf{1.} $1\,\text{نانومول}/5\,\text{لتر} = \qty{0.2}{nmol/L}$.
\textbf{2.} $\theta = 0.2/1.2 = 0.17$: أي نحو 3300 مستقبل مشغول.
\textbf{3.} ثلاثة أعمار نصف: \qty{0.025}{nmol/L}؛ و$\theta = 0.024$.
\textbf{4.} نحو نصف دقيقة إلى دقيقة، وهو زمن دورة
الدم؛ أما العصب فيسلّم ناقله مباشرةً إلى
الكبد في غضون ثانية.
\textbf{5.} الكسر نفسه، لكن $33\,000$ مستقبل مشغول: أي
استجابة أكبر بعشرة أضعاف.
\textbf{6.} عند \qty{1}{nmol/L} يكون مستقبل $K_d = \qty{1}{\micro
mol/L}$ مشغولًا بمقدار جزء من ألف: فلا استجابة. ومطابقة $K_d$
للمجال الدائر تضع الجزء شديد الميل من منحني الإشغال حيث
يتغير \omterm{def:b2:cell-signalling:kinds}{الهرمون} فعلًا.
\textbf{7.} $20\times 100\times 1\times 50\times 50\times 1000 =
5\times 10^{9}$ جزيء غلوكوز في الثانية لكل مستقبل مشغول.
\textbf{8.} $3300\times 20\times 100 = 6.6\times 10^{6}$ من \omterm{prop:b2:cell-signalling:gpcr}{AMP الحلقي} في
الثانية الأولى؛ وفي \qty{5e-12}{L} يكون ذلك $1.1\times 10^{-17}$ mol، أي نحو
\qty{2}{\micro mol/L}.
\textbf{9.} $3300\times 5\times 10^{9} = 1.7\times 10^{13}$ في الثانية، و$10^{15}$
في الدقيقة للخلية الواحدة.
\textbf{10.} $2\times 10^{11}$ خلية $\times 10^{15} = 2\times
10^{26}$ جزيء في الدقيقة، أي 330 مولًا، و\qty{6e4}{g} في الدقيقة --- أي عشرة
آلاف ضعف \qty{5}{g/min} الملاحظة. فالربوح لا تتضاعف
إلا ما دامت كل مرحلة غير مشبعة؛ والفسفوريلاز في الواقع
محدود بركيزته وبالفوسفاتازات، وتصدير
الغلوكوز محدود بالناقلات. فربح الشلال يُنفق على السرعة
والحساسية، لا على الخرج.
\textbf{11.} عند \qty{5}{g/min}، عشرون دقيقة؛ والفزع ينقضي في
دقائق ومفاتيح الإطفاء توقف الشلال قبل ذلك بكثير.
\textbf{12.} حلمهة GTP \omterm{prop:b2:cell-signalling:gpcr}{ببروتين G} (ثوانٍ)؛
والفوسفودي إستيراز على \omterm{prop:b2:cell-signalling:gpcr}{AMP الحلقي} (ثوانٍ)؛ والفوسفاتازات على
الإنزيمات المفسفرة (من ثوانٍ إلى دقائق)؛ وفقد حساسية
المستقبل (دقائق). وأسرعها ساعة \omterm{prop:b2:cell-signalling:gpcr}{بروتين G} نفسها
والفوسفودي إستيراز.
\textbf{13.} لا يُتلف \omterm{prop:b2:cell-signalling:gpcr}{AMP الحلقي}: فيتراكم ويدوم، ويبقى
الكيناز مشتغلًا، ويكون خرج الغلوكوز أكبر وأطول
--- وهي يقظة القهوة.
\textbf{14.} $4000/50 = \qty{80}{s}$.
\textbf{15.} $600/5 = \qty{120}{s}$ لكل إنزيم لكل ريباس؛ أي 30 لكل
ريباس في الساعة، و600 لكل mRNA في الساعة.
\textbf{16.} $100\times 600 = 60\,000$ إنزيم في الساعة؛ ويُبلغ $10^{5}$ بعد
\qty{1.7}{h} نحوًا.
\textbf{17.} من عشر دقائق إلى عشرين ليدخل الكورتيزول ويرتبط
ويبلغ مورّثاته، وساعة قبل أن تكثر أول نسخ mRNA وتُترجم،
ثم ساعتان من التراكم: أي نحو أربع ساعات.
\textbf{18.} الشلال يغير فعالية إنزيمات موجودة؛
وعمل الكورتيزول أن يغير أي الإنزيمات يوجد، لأيام، وهو ما
لا يقدر عليه إلا الاستنساخ. والفزع لا يستطيع انتظار أربع ساعات لإنزيمات
جديدة.
\textbf{19.} مستقبلات أكثر تعني مستقبلات مشغولة أكثر عند أي
تركيز للأدرينالين: فكسر الإشغال لا يتغير لكن
الاستجابة عند كل تركيز أكبر --- فالمنحني
مضروب في عامل، لا مزاح. \omterm{def:b2:cell-signalling:kinds}{فالهرمون} البطيء يضبط ربح السريع.
\textbf{20.} يدخل الغلوكوز الخلية $\beta$ ويُستقلب؛ فيرتفع ATP
ويغلق قناة بوتاسيوم حساسة لمستوى ATP؛ ويستقطب الغشاء؛
وتنفتح قنوات كالسيوم فولطية البوابة؛ ويطلق الكالسيوم
إخراج حبيبات الأنسولين --- أي استقلاب مقترن
بالإفراز بخطوة كهربائية.
\textbf{21.} يبدأ المستقبلان طريقين يلتقيان على
الإنزيمات نفسها: فكيناز A التابع للغلوكاغون يفسفرها، وشلال الأنسولين
يفعّل الفوسفاتاز التي تنزع فوسفاتها؛ وحالة
الإنزيمات، ومن ثم اتجاه استقلاب الغليكوجين، تتبع
نسبة الاثنين.
\textbf{22.} الغلوكوز الصائم مرتفع (فالكبد، بلا كابح، يمضي
في صنعه)، والغلوكوز بعد الوجبة مرتفع جدًا (فلا التقاط في العضل
والشحم)؛ ودون الأنسولين تطلق مخازن الشحم حموضًا دسمة، يحولها
الكبد إلى حموض كيتونية --- وهو الحماض الكيتوني.
\textbf{23.} تلتقط العضلة العاملة الغلوكوز دون أنسولين، فيهبط
الغلوكوز في عَدْو طويل بينما يمضي الأنسولين المحقون، بلا تنظيم،
في كبح الكبد: فنقص سكر الدم والانهيار.
\textbf{24.} حلقة الغلوكوز: المستشعر هو الخليتان $\beta$ و$\alpha$،
وهرمونان منفِّذان للاتجاهين، وربح عالٍ يمسك
المستوى إلى حدود الخُمس، ودقائق للفعل. أما \omterm{def:b2:blood-pressure:baroreflex}{المنعكس الضغطي}: فمستقبلات
التمدد، \omterm{def:b2:heart:anatomy}{والقلب} والأوعية منفِّذات، وربح نحو أربعة، و
ثانية للفعل.
\textbf{25.} الإشغال $0.17$ بعد الإطلاق؛ والربح $5\times 10^{9}$ في الثانية
لكل مستقبل مشغول؛ وخرج الكبد \qty{5}{g/min} في الواقع؛ والكورتيزول
يستغرق نحو أربع ساعات.
\end{solution}
